\chapter{브레이크}
% full version: chapters/03_gaba.tex

기계에 수가 두 번째로 많은 뉴런 유형인 \nt{GABA}를 더하자. 이들의 신호는 이웃을 딸깍 쪽으로 밀지 않고 밀어낸다. 이것이 브레이크이다. 수는 많지 않아서 피질에서도 전체 뉴런의 5분의 1에서 4분의 1 정도이지만, 이들이 없으면 기계는 굴러가지 않는다.

제약이 하나 있다. “플러스” 뉴런은 이웃을 직접 억제할 수 없다. 음의 연결이 필요하면 중개자를 거쳐야 한다. “플러스”가 억제 뉴런을 깨우고, 그 억제 뉴런이 목표를 억제한다. 부호를 바꾸는 값은 뉴런 하나와 작은 지연 하나이다.

\section*{“단, B가 오면 빼고”}

브레이크로 할 수 있는 가장 쓸모 있는 연산은 금지이다. “A가 오면 발화하라, 단 B가 왔으면 말고.” 뉴런은 A에게서 밀어 주는 신호를, B가 깨우는 중개자에게서 브레이크를 받는다. 금지와 “또는”만 있으면 어떤 논리든 조립할 수 있다.

\pencilfig{3}{\pencilart{3}}{“A, 단 B가 오면 빼고”: 빨간 뉴런이 파란 뉴런의 길을 막는다.}

시간 속의 금지는 여러분 눈에 있는 운동 검출기가 된다. 이웃한 감지기 두 개가 있다. 하나는 출력 뉴런을 흥분시키고, 다른 하나는 지연을 두고 억제한다. 빛 점이 알맞은 속도로 한쪽 방향으로 달리면 억제가 흥분과 정확히 같은 때 도착해 흥분을 꺼 버린다. 반대 방향이면 억제가 늦고, 뉴런은 그전에 딸깍한다. 한 방향의 움직임만 보는 뉴런이 생긴다.

\pencilfig{103}{%
% sensors on top; the left one excites the output, the right one inhibits it through a GABA go-between and a delay
\draw[dotpen=pcAmber, wash=pcAmber] (0,1.6) circle (0.16); \draw[dotpen=pcAmber, wash=pcAmber] (1.6,1.6) circle (0.16);
\node[lbl, anchor=south] at (0.8,1.8) {감지기 두 개};
\draw[pen=pcBlue, wash=pcBlue] (0.4,0) circle (0.24);
\draw[arr=pcBlue] (0,1.44) -- (0.3,0.25);
\draw[pen=pcBlue] (1.6,1.44) -- (1.6,1.18);
\draw[penlite, fill=white] (0.95,0.9) rectangle (2.25,1.18); \node[font=\tiny\itshape, text=pcGray] at (1.6,1.04) {지연};
\draw[pen=pcBlue] (1.6,0.9) -- (1.6,0.72);
\draw[pen=pcRed, wash=pcRed] (1.6,0.55) circle (0.17);
\draw[pcRed, line width=0.7pt, -{Bar[width=5pt]}] (1.45,0.45) -- (0.66,0.08);
\draw[arr] (0.4,-0.24) -- (0.4,-0.6); \node[lbl, anchor=north] at (0.4,-0.6) {출력};
% outcomes
\begin{scope}[xshift=3.1cm]
  \draw[dotpen=pcAmber, wash=pcAmber] (0,1.25) circle (0.1); \draw[arr=pcAmber] (0.18,1.25) -- (1.1,1.25);
  \node[lbl, anchor=west] at (1.25,1.25) {왼쪽→오른쪽: 딸깍};
  \draw[dotpen=pcAmber, wash=pcAmber] (1.1,0.45) circle (0.1); \draw[arr=pcAmber] (0.92,0.45) -- (0,0.45);
  \node[lbl, anchor=west] at (1.25,0.45) {오른쪽→왼쪽: 침묵};
  \node[lbl, anchor=west, align=left] at (0,-0.35) {억제가 흥분과\\정확히 같이 도착};
\end{scope}
}{운동 검출기: 억제는 지연되어 도착하므로, 빛 점이 오른쪽에서 왼쪽으로 달릴 때만 응답을 끈다.}


\section*{뺄셈인가 나눗셈인가}

억제는 빼기만 하는 것이 아니다. 열린 억제 채널은 전압을 아래로 끌어내리는 것이 아니라 \emph{안정 상태 쪽으로} 끌어당긴다. 양동이에 구멍을 하나 더 뚫은 것과 같다. 그러면 어떤 물이 들어와도 수위가 덜 오른다. 억제는 신호에서 일정한 몫을 덜어 내는 것이 아니라 신호를 \emph{나눈다}. 볼륨 조절기이다. 억제의 세기가 주변의 전체 활동과 함께 커지면, 각 뉴런은 자기 신호가 얼마나 센지가 아니라 \emph{이웃에 비해} 얼마나 센지에 반응한다. 그래서 같은 네트워크가 어스름에서도 한낮의 햇빛 아래서도 일한다. 다만 이런 “나눗셈기”가 스파이크 빈도에 깔끔하게 작용하려면 일정한 배경 잡음이 함께 있어야 한다. 살아 있는 뇌에는 그런 잡음이 늘 있다.

\pencilfig{104}{%
\foreach \dx/\t in {0/뺄셈,4.6/나눗셈}{
  \begin{scope}[xshift=\dx cm]
  \draw[pen] (0,0) -- (3.6,0); \draw[pen] (0,0) -- (0,1.9);
  \node[lbl, anchor=north] at (1.8,-0.05) {신호}; \node[lbl, anchor=south, rotate=90] at (-0.05,0.95) {응답};
  \draw[pen=pcBlue] (0.4,0) -- (3.3,1.75);
  \node[font=\small\itshape, text=pcGray, anchor=south] at (1.8,1.95) {\t};
  \end{scope}}
\draw[pen=pcRed, line width=0.9pt] (1.3,0) -- (3.4,1.27);
\node[lbl, text=pcRed, anchor=west] at (2.25,0.35) {이동};
\begin{scope}[xshift=4.6cm]
\draw[pen=pcRed, line width=0.9pt] (0.4,0) -- (3.3,0.8);
\node[lbl, text=pcRed, anchor=west] at (2.0,0.25) {기울기 감소};
\end{scope}
\node[lbl, text=pcBlue, anchor=east] at (2.25,1.45) {억제 없음};
}{파란색은 억제가 없을 때, 빨간색은 억제가 있을 때 뉴런의 응답이다. 뺄셈은 문턱값을 옮기고, 나눗셈은 볼륨 노브처럼 소리를 줄인다.}

두 번째 결과도 있다. 구멍이 하나 더 있으면 양동이가 더 빨리 잊는다. 두 신호가 더해지려면 들어와야 하는 창이 좁아진다. 억제는 동시성 검출기를 더 엄격하게 만든다.

\section*{선택}

이제 모자랐던 가장 중요한 것, 여럿 중 하나를 고르는 일이다. 두 뉴런 무리가 서로를 억제한다. 억제가 약한 동안에는 둘 다 일하되 더 조용히 일한다. 억제가 충분히 세면 작은 차이도 그네처럼 점점 커져서, 결국 한 무리가 완전히 입을 다문다. 승자독식이다. 그리고 승리는 유지된다. 네트워크가 마음을 바꾸려면 새 근거가 옛 근거보다 눈에 띄게 강해야 한다. 기계는 바스락 소리마다 움찔하지 않는다.

브레이크는 기억도 지운다. 억제 뉴런이 “리셋” 신호를 받아 지난 장의 모닥불을 끈다. 서로를 억제하는 이런 셀 두 개는 래치에 해당한다. 모든 컴퓨터 메모리의 기초이다.

\section*{경계 위의 삶}

“플러스”로만 된 네트워크는 사태로 치달을 수 있다. 스파이크 하나가 다음 스파이크를 하나보다 많이 낳는 상태이다. 브레이크는 오븐의 온도 조절기가 온도를 붙잡듯 네트워크를 동작점에 붙잡아 둔다. 조건은 둘이다. 억제는 충분히 세고 충분히 빨라야 한다. 세지만 느리면 조절기가 지나치게 반응한다. 샤워할 때 뜨거운 물이 올 때까지 기다리지 못하고 수도꼭지를 돌려 대는 사람과 같다. 온도는 이리저리 흔들린다. 뇌에서는 이런 “플러스”와 “마이너스”의 고리가 빠른 리듬, 즉 1초에 수십 번의 진동을 만든다.

피질은 바로 이런 경계 위에서 사는 것으로 보인다. 피질 자신의 되먹임은 브레이크가 없으면 사태로 치달을 만큼 세다. 이를 붙잡는 것은 빠른 브레이크뿐이며, 그래서 피질은 약한 신호에 그토록 예민하게 반응한다.

이 장의 결론은 뜻밖이다. 브레이크는 기계에 새로운 계산을 거의 더하지 않는다. 브레이크가 더하는 것은 \emph{제어}이다. 선택, 리셋, 볼륨 조절, 안정성. 흥분은 데이터를 나르고, 억제는 어떤 데이터가 언제, 얼마나 크게 지나갈지 정한다.

\fullversion{3}
